% secctions/1_introduccion.tex
\section{Introducción a \LaTeX\ y su Arquitectura Documental}

\subsection{Origen y Definición de \TeX\ y \LaTeX}
\TeX\ es un sistema de composición de documentos de alta calidad, orientado especialmente a la creación de documentos científicos y técnicos que incluyen fórmulas matemáticas. Este sistema fue creado por el científico de la computación Donald Knuth en 1978. A diferencia de un procesador de textos de uso general, \TeX\ no es una aplicación interactiva, sino un lenguaje de programación que requiere compilar el código fuente para obtener el documento final.

Para abstraer la complejidad nativa de este sistema, Leslie Lamport desarrolló \LaTeX, el cual es un conjunto de macros para \TeX\ diseñadas originalmente para facilitar su uso. En el ámbito académico y de desarrollo, es fundamental destacar que tanto \TeX\ como \LaTeX\ son programas de código abierto, liberados bajo la licencia LPPL.

\subsection{El Paradigma de Compilación frente a WYSIWYG}
La diferencia fundamental entre \LaTeX\ y los procesadores de texto comerciales, como Microsoft Word, reside en su filosofía de diseño. Los procesadores de texto tradicionales siguen el paradigma WYSIWYG (\textit{What You See Is What You Get}). 

En contraposición, el enfoque de compilación de \LaTeX\ es, en realidad, su mayor ventaja técnica, ya que permite separar fácilmente el contenido y la estructura de un documento, de su formato. Esta arquitectura permite que el usuario pueda centrarse exclusivamente en el contenido y la estructura del documento, delegando a \TeX\ la tarea de encargarse de aplicar un diseño tipográfico estandarizado y riguroso.

\subsection{Estructura Fundamental de un Documento}
El código fuente de cualquier proyecto en \LaTeX\ requiere una estructura jerárquica estrictamente definida. El esqueleto básico se divide en dos bloques principales:

\begin{itemize}
    \item \textbf{La Clase y el Preámbulo:} La primera línea de un fichero con código \LaTeX\ indica la clase de documento que se va a generar mediante el comando \texttt{\textbackslash documentclass}. El preámbulo es la parte que va después de la clase y antes del comienzo del cuerpo del documento. Suele utilizarse para la carga de los paquetes de macros y la configuración global. Aquí se configuran elementos críticos como el paquete \texttt{inputenc} para la codificación de caracteres, \texttt{babel} para definir el idioma del documento, y \texttt{fontenc} para las codificaciones de las fuentes.
    
    \item \textbf{El Cuerpo del Documento:} Contiene el texto del cuerpo del documento y tiene que ir dentro del entorno \texttt{document}. Todo el contenido visible debe encapsularse estrictamente entre las directivas \texttt{\textbackslash begin\{document\}} y \texttt{\textbackslash end\{document\}}. Dentro de esta sección suele empezar con el comando \texttt{\textbackslash maketitle} para generar la portada y le sigue \texttt{\textbackslash tableofcontents} que introduce la tabla de contenidos automáticamente.
\end{itemize}